\chapter{Dry Mouth}

\section*{Definition}
Reduced salivary flow causing oral mucosa dryness difficulty swallowing tasting speaking altering normal mouth sensation.

\section*{Pathophysiology}
Xerostomia can lead dental caries periodontal disease taste disturbances oral infections halitosis speaking eating difficulties; distinguished from normal temporary dryness persists despite hydration medication induced systemic disease radiation therapy Sjogren syndrome autoimmune diabetes aging dehydration chemotherapy opioids antihistamines antidepressants.

\section*{Causes}
\begin{itemize}
  \item Medications (anticholinergics antihistamines)
  \item Autoimmune disorders (Sjogren syndrome)
  \item Radiation therapy head/neck
  \item Diabetes mellitus
  \item Dehydration/insufficient fluid intake
\end{itemize}

\section*{When to See a Doctor}
See your doctor if:
\begin{itemize}
  \item Severe or worsening symptoms
  \item Symptoms lasting more than persistent beyond simple dehydration unexplained accompanied other systemic symptoms warrants dental medical evaluation
  \item Accompanied by other concerning signs
\end{itemize}

\section*{Related Symptoms}
\begin{itemize}
  \item Difficulty swallowing dry throat tooth decay bad breath altered taste frequent thirst
\end{itemize}
\textbf{Important:} If this symptom persists, your doctor will check for serious underlying causes.


\section*{Self-Care / Home Management}
\begin{itemize}
  \item Maintain good oral hygiene: brush twice daily, floss daily, and rinse with warm salt water (1/2 tsp salt in 8 oz water).
  \item Avoid tobacco, alcohol, and spicy or acidic foods that may irritate oral tissues.
  \item Stay hydrated by sipping water throughout the day to combat dry mouth.
  \item Over-the-counter pain relievers and topical oral gels can provide temporary relief for toothache or gum discomfort.
  \item See a dentist promptly for persistent oral symptoms, as early intervention improves outcomes.
\end{itemize}

\section*{How Your Doctor Will Evaluate You}
\begin{itemize}
  \item History covers onset, character of pain, aggravating/relieving factors, oral hygiene habits, and dental care history.
  \item Oral examination inspects teeth, gums, mucosa, tongue, and oropharynx for caries, inflammation, lesions, or masses.
  \item Dental X-rays (panoramic, periapical) identify caries, abscesses, impactions, and bone pathology.
  \item Referral to a dentist or oral surgeon is appropriate for definitive diagnosis and treatment of dental pathology.
\end{itemize}

\section*{Treatment Options}
\begin{itemize}
  \item Dental caries require restorative treatment (fillings, crowns, root canal) by a dentist.
  \item Oral infections (abscess, pericoronitis) may require drainage, antibiotics, and definitive dental treatment.
  \item Dry mouth management includes hydration, sugar-free lozenges, saliva substitutes, and addressing underlying causes.
  \item Regular dental check-ups every 6 months are essential for prevention and early detection of oral pathology.
\end{itemize}

\clearpage
